% Part: incompleteness
% Chapter: theories-computability
% Section: complete-decidable

\documentclass[../../../include/open-logic-section]{subfiles}

\begin{document}

\olfileid{inc}{tcp}{cdc}

\olsection{\printtoken{S}{axiomatizable} संपूर्ण उपपत्ती निर्णेय असतात}

प्रत्येक वाक्य~$!A$ साठी $!A$ किंवा $\lnot !A$ यांपैकी
एक सिद्ध होत असेल, तर उपपत्तीला {\em संपूर्ण} म्हणतात.

\begin{lem}
उपपत्ती $\Th{T}$ संपूर्ण आणि !!{axiomatizable} आहे असे समजू.
मग $\Th{T}$ निर्णेय असते.
\end{lem}

\begin{proof}
$\Th{T}$ संपूर्ण आहे आणि $A$ हा तिचा संगणनक्षम स्वयंसिद्धकसंच
आहे असे समजू. $\Th{T}$ विसंगत असेल, तर ती स्पष्टपणे संगणनक्षम
आहे. (अल्गोरिदम: ``नेहमी होय म्हणा.'') म्हणून $\Th{T}$
सुसंगतही आहे असे मानू शकतो.

वाक्य $!A$ हे $\Th{T}$ मध्ये आहे की नाही हे ठरवण्यासाठी,
$\Th{T}$ पासून $!A$ ची !!a{derivation} आणि
$\lnot !A$ ची !!a{derivation} यांचा एकाच वेळी शोध घ्या.
$\Th{T}$ संपूर्ण असल्यामुळे यांपैकी एक तरी निष्पत्ती मिळेलच;
आणि $\Th{T}$ सुसंगत असल्यामुळे $\lnot !A$ ची
!!a{derivation} मिळाल्यास $!A$ ची !!{derivation} असणार नाही.

हेच दुसऱ्या प्रकारे पाहू: $\Th{T}$ ही !!{c.e.} आहे, हे आपण
आधीच जाणतो. म्हणून आधी सिद्ध केलेल्या प्रमेयानुसार,
$\Th{T}$ चा पूरक संचदेखील !!{c.e.} आहे हे दाखवणे पुरेसे आहे.
पण !!a{formula}~$!A$ हे $\Th{\bar
T}$ मध्ये असते तेव्हाच आणि तरच $\lnot !A$ हे $\Th{T}$
मध्ये असते. त्यामुळे $\Th{\bar T} \leq_m
\Th{T}$.
\end{proof}

\end{document}
